\section{Introduction}\label{intro:section}

Consider an LP vertex that violates a quadratic constraint. A basis gives
a cone with that vertex as its apex. A closed convex set whose interior
contains the apex and avoids the quadratic feasible set gives an
intersection cut: follow each cone ray to the boundary of the convex set,
then cut off the simplex formed by those ray intersections. Larger ray
steps give smaller cut coefficients. The choice of the convex set
therefore matters, even when the vertex, basis, and violated constraint
are fixed. Intersection cuts for integer programming go back to Balas
\citep{Balas1971}.

Quadratic constraints permit many such sets. A set selected from the
violated point need not give the best objective bound on the corner. A
set giving that bound need not yield the strongest bound after the LP is
reoptimized. A family whose best single cut is weak may improve when all
its cuts are imposed, but even that closure can remain weaker than the
nonlinear corner. Finally, good cuts at individual vertices do not by
themselves imply convergence of successive LP relaxations. These are
different questions, and answering one does not settle the others.

This paper develops them in a common framework. Our principal examples
are the bilinear inequality $w-xy\le0$ and the homogeneous determinant
equation $ad-bc=0$. Both arise in lifted formulations of quadratic
optimization. The bilinear inequality also isolates a difficulty that is
easy to hide in a larger model: its maximal free sets can extend far
beyond the part of the corner that determines the objective bound.
Studying this geometry gives exact limits on restricted set selection,
quantitative guarantees when the corner is suitably scaled, and
conditions for repeated cuts to converge.

\subsection{Prior work and the question studied here}

The relation between valid inequalities and free sets is established in
the cut-generating-function literature. A convex free neighborhood of
the origin defines a gauge and hence an inequality for a corner
relaxation. \citet{CCDLM2015} develop this correspondence and explain why
not every valid inequality is generated by an ordinary, finite-valued
cut-generating function. Their Theorem~6.3 gives sufficiency when the
projected rays positively span the ambient space.
\citet{CWY2015}, Theorem~1.1, prove exact domination of every valid
inequality under the more general condition that the translated feasible
set lies in the ray cone. \citet{KKSteffy2016}, Proposition~5, recover
this result, while their Corollary~2 proves sufficiency of relaxed
cut-generating functions without that cone condition. A relaxed function
can take the value $+\infty$, and its level set need not be a neighborhood
of the origin.

These distinctions are relevant to the unrestricted benchmark in this
paper. Equality of the best intersection-cut bound and the corner bound
is a supremum statement. An ordinary free neighborhood need not attain
that supremum. \citet{KKYang2017} study sufficient
and necessary conditions for ordinary cut-generating functions outside
the cone-containment setting. Its Examples~4.2 and~4.3 illustrate the
difference between transverse and tangential approach to a boundary
point. Our general transversality criterion gives a sufficient condition
for differentiable sublevel sets, stated at every minimizing contact of
the corner. It applies without the cone-containment condition, and its
proof uses the geometry of these contacts directly. The unrestricted
exactness and dominance results serve as background; we do not claim
them as new.

Choosing deeper cuts also has a substantial history.
\citet{EN2005} formulate a distance-based deepest convexity-cut problem
over structured free sets for a finite mixed integer point set.
\citet{PoirrierYu2019} bound Euclidean cut depth for
tableau corners and identify set-dependent depth as a possible criterion
for selecting lattice-free sets. \citet{XavierFukasawaPoirrier2021} minimize
the largest coefficient in a multi-row lattice-free construction. Our
one-cut bound is the reciprocal of the largest normalized coefficient
$a_j(C)/\omega_j$. These objectives coincide when the reduced costs are
equal; with unequal costs, ours is a weighted version. Coefficient
optimization therefore has a direct precedent.
\citet{BalasMargot2013} study a
complementary choice: a fixed free set with a tighter polyhedral
relaxation than a basis cone. Thus optimizing a free set, or seeking a
deeper intersection cut, is not a new general idea. Our selection
criterion is the objective bound on a fixed corner, and our restricted
families exclude a continuous quadratic feasible set. Euclidean depth
enters separately when we study successive rounds.

Family closures have also been studied extensively for lattice-free
sets. \citet{ABP2018}, Theorem~5, relate constant-factor approximation
by a family closure to approximation of each target cut by a single
family member, under their closedness and representation hypotheses.
That result concerns lattice-free corner polyhedra at a fixed apex.
Our closure statements concern continuous quadratic feasible sets and
specific determinant families. Their geometric counterexamples and
factor classifications do not follow by transferring the lattice-free
theorem. The common point is that combining cuts and approximating a
target closure must be distinguished from obtaining one exact cut.

For quadratic constraints, \citet{MPS2022} construct and analyze
quadratic-free sets used to generate intersection cuts.
\citet{MPS2025} characterize maximal free sets for homogeneous quadratic
inequalities, and \citet{MPS2026} treat the inhomogeneous case.
These characterizations describe the available sets. They do not imply
that a set chosen from the violated point optimizes a prescribed corner
objective. Our low-dimensional classification is a consequence of this
geometry and Lorentz transformations, rather than a replacement for the
general characterization.

Outer-product formulations give another source of quadratic-free sets.
\citet{BCM2020}, Lemma~22 and Theorem~23(iv), construct maximal
outer-product-free sets for $2\times2$ minors. Their rotation family is
contained in the determinant orbit studied here. Their discussion also
explicitly notes that a choice that is best in a violation sense need not
produce the deepest cut. We enlarge that rotation family to a
three-parameter orbit, distinguish it from the transformed point rule,
and optimize its corner bound by small semidefinite feasibility problems.
The existence of the rotation subfamily is prior work.

\citet{CMS2023} implement quadratic and minor intersection cuts in SCIP.
The experiments in their earlier report \citep{ChmielaMunozSerrano2020ZIB}
concern root bounds; the best reported configuration combines quadratic
intersection cuts with minor cuts. A separate bilinear implementation studied by
\citet{FischettiMonaci2020} has heterogeneous computational effects,
including improvements on selected instances and losses in its overall
comparison. These studies motivate careful separation of local cut
strength, root relaxation strength, and branch-and-bound work. Our
computational analysis makes that separation explicit and reports
implementation and numerical qualifications that affect the conclusions.

\subsection{Results and contributions}

Let $\sbar$ be the projected corner vertex, let $p_j$ be its projected
rays, and let $\omega_j>0$ be their reduced costs. The nonlinear corner
bound is
\[
 \zK(\rc)=\inf\{\rc^T\lambda:
       \lambda\ge0,\ \sbar+\textstyle\sum_j\lambda_jp_j\in S\}.
\]
If a free set $C$ has ray steps $\alpha_j(C)$, its one-cut bound is
$\min_j\omega_j\alpha_j(C)$, with infinite steps handled as specified in
\cref{fd:single-cut}. This formula turns a geometric set-selection
problem into an objective-specific one.

Our main results concern the restrictions that remain after this
unrestricted benchmark is understood.

\begin{enumerate}
\item \textbf{Bilinear orbit bounds and their limits.}
For $M(x,y,w)=\begin{psmallmatrix}w&x\\y&1\end{psmallmatrix}$, the sets
$C_F=\{s:\sym(F^TM(s))\succeq0\}$, $\det F>0$, form a
three-parameter family. We study this family $\famA$ and its closed
upward completions $\famB$. Rational corners give strict gaps for both
families, although an unrestricted free set attains the corner bound.
Tangent edges force the orbit parameter onto a pencil, which explains
several of these obstructions. A support-one example, whose minimizing
ray vector has exactly one positive component, shows that the obstruction
is not confined to tangent edges.

\item \textbf{Quantitative guarantees.}
Scale each ray by the distance corresponding to the true corner bound.
A bilinear measure $D$ of the resulting ray size relative to vertex
violation gives a positive lower bound for $\zA/\zK$ and $\zB/\zK$.
Its worst-case order $1/D$ is sharp
(\cref{dp:depth-bound,dp:vanishing,dp:certified-sharpness}).
For SCIP's fixed point rule in a specified quadratic representation, the
sharp order is $\min\{1,1/(\kappa D)\}$, where $\kappa$ measures
conditioning of the scaled three-ray matrix
(\cref{dp:fixed-rule-bound,dp:fixed-rule-sharp}).
These statements have different invariance properties; the point-rule
bound is not a representation-independent solver guarantee.

\item \textbf{Contact, attainment, and closures.}
At a support-one bilinear contact, interval conditions characterize
orbit supremum equality and distinguish it from attainment
(\cref{dp:contact-intervals}). Taking every cut from a restricted family
can improve a single-cut bound without recovering the true corner
bound. Exact certificates establish strict fixed-corner closure gaps.
Upward completion can change an infinite approximation factor into
exactness at one corner, yet neither $\famA$ nor $\famB$ has a uniform
positive approximation factor over all corners.

\item \textbf{Homogeneous minors.}
For $ad-bc=0$, the same determinant orbit consists of maximal free cones.
The transformed point rule has two parameters; the rotation family has
one, and their intersection at a given vertex is its polar-factor set.
Full-dimensional rational examples prove strict orbit gaps, including a
support-one example. Complete rational matrix witnesses certify the
headline gaps (\cref{mi:examples}).

\item \textbf{Successive cuts and computation.}
A uniform positive cut depth at vertices with a fixed positive violation
ensures convergence on a compact relaxation. Uniform pointedness of
the basis cones supplies a sufficient geometric condition for specified
rules. An explicit infinite sequence of maximal bilinear-free sets has
uniformly pointed basis cones and fixed violation, but its LP bounds
converge below the nonlinear optimum. Archived experiments then compare
the different strength criteria. Better corner-bound selection does not
consistently yield better reoptimized LP bounds, fewer rounds, or better
full-solve performance.
\end{enumerate}

To the best of our knowledge, the particular quantitative guarantees,
contact criteria, rational restricted-family and closure obstructions,
and explicit nonconvergent sequence stated here are not contained in the
prior work reviewed above. This claim concerns these precise statements,
not the general use of quadratic-free sets, determinant transformations,
free-set optimization, or cutting-plane convergence arguments. The
rank and inertia bounds, corner algorithms, and low-dimensional geometry
also provide tools for deriving and checking the examples. We identify
their established ingredients where they are used.

\subsection{How the results fit together}

The comparison is always made at a specified level. The following table
records the objects that must be distinguished.

\begin{center}
\begin{tabular}{p{0.27\textwidth}p{0.64\textwidth}}
\toprule
Object & Question answered \\
\midrule
Nonlinear corner bound $\zK$ & What can the violated constraint enforce
on the chosen corner? \\
Best single family cut & How much of that bound can one selected set
enforce? \\
Fixed-corner family closure & How much can all sets in that family enforce
on the same corner? \\
LP reoptimization & What remains after imposing a cut together with all
other LP rows? \\
Successive rounds & Do new corners and cuts approach the nonlinear
optimum? \\
Branch-and-bound & How do separation cost, numerical behavior, and the
changed search tree affect the solve? \\
\bottomrule
\end{tabular}
\end{center}

\Cref{fd:section} establishes the corner benchmark and computational
reductions. \Cref{fd:geometry-section} introduces the quadratic geometry
and bilinear orbit. \Cref{dp:section} proves the quantitative and contact
results, and \cref{cl:section} treats fixed-corner closures.
\Cref{mi:section} gives the homogeneous-minor extension.
\Cref{cv:section} separates sufficient convergence conditions from
counterexamples. \Cref{exp:section} presents the archived computational
evidence. The appendices contain complete technical proofs and the
certificate specification. The electronic companion retains the exact
certificate data and the compact experimental records.
